\vfill%
\pagebreak

\section*{Exercises}
\markright{\MakeUppercase{Exercises}}%
\subsection*{Exercise 1}

Write a Ymir program that declares a mutable \token{i32} variable named
\token{score}, initialized to \token{5}, and prints it. Then, update
\token{score} so that it holds its own square, and print it again.

\subsection*{Exercise 2}

What does the following program print? Work out the result by hand using the
precedence rules from Table~\ref{tab:(chap2):operator_precedence}, then check
your answer by compiling and running the program.

\begin{lstlisting}[style=coloredverbatim]
use std::io;

fn main() {
  let a = 2;
  let b = 5;
  let c = 3;

  let r = a + b * c - a;
  println("R : ", r);
}
\end{lstlisting}

\subsection*{Exercise 3}

The following program contains two errors, one for each kind of mistake this
chapter introduced.

\begin{lstlisting}[style=coloredverbatimError, escapechar=@, caption=Ymir program with errors]
use std::io;

fn main() {
  let x: i32 = 10;
  @\hb{x = 20;}@

  let y: f32 = @\hb{"text"}@;

  println("X : ", x);
  println("Y : ", y);
}
\end{lstlisting}

Correct both errors and compile the program with \token{gyc} to test your
corrections.

\subsection*{Exercise 4}

What does the following program print? \token{u8} is an unsigned 8-bit
integer type.

\begin{lstlisting}[style=coloredverbatim]
use std::io;

fn main() {
  let mut x: u8 = 250;

  x += 10;
  println("X : ", x);
}
\end{lstlisting}

\subsection*{Exercise 5}

What does the following program print? For each of the five variables, state
which floating-point or integer type it holds, and explain why.

\begin{lstlisting}[style=coloredverbatim]
use std::io;

fn main() {
  let a = 12;
  let b = 12u64;
  let c = 3.14;
  let d = 3.14f;
  let e = 1_000_000;

  println("A : ", a);
  println("B : ", b);
  println("C : ", c);
  println("D : ", d);
  println("E : ", e);
}
\end{lstlisting}

\subsection*{Exercise 6}

What does the following program print? The word \textit{café} has four
characters. Explain why the printed value is not 4.

\begin{lstlisting}[style=coloredverbatim]
use std::io;

fn main() {
  let word = "café";

  println(word.len);
}
\end{lstlisting}

\subsection*{Exercise 7}

What does the following program print? List, in order, which calls to
\token{check} are actually made, and explain why the other ones are skipped.

\begin{lstlisting}[style=coloredverbatim]
use std::io;

fn check(name: [c8], value: bool)-> bool {
  println("Checking ", name);
  value
}

fn main() {
  let mut ok = true;

  if check("a", false) && check("b", true) {
    ok = false;
  }

  println("OK : ", ok);
}
\end{lstlisting}

\vfill%
\pagebreak

\forceevenpage{}

\section*{Solutions}
\markright{\MakeUppercase{Solutions}}%
\subsection*{Exercise 1}

\begin{lstlisting}[style=coloredverbatim, caption=Solution for exercise 1]
use std::io;

fn main() {
  let mut score = 5;
  println("Score : ", score);

  score = score * score;
  println("Score : ", score);
}
\end{lstlisting}

\subsection*{Exercise 2}

The multiplication \token{b * c} is evaluated first, since \token{*} has
higher precedence than \token{+} and \token{-}. The remaining two operators
share the same precedence, so they are evaluated left to right:
\token{(a + (b * c)) - a}, that is \token{(2 + 15) - 2 = 15}.

\begin{lstlisting}[style=bashVerb]
R : 15
\end{lstlisting}

\subsection*{Exercise 3}

The variable \token{x} is declared with \token{let} but not \token{let mut},
so it is a constant and cannot be reassigned. The variable \token{y} is
declared with the type \token{f32}, but the value on the right-hand side of
the \token{=} sign is a string literal, not a floating-point literal — the
types of the \textit{lvalue} and the \textit{rvalue} do not match.

\begin{lstlisting}[style=coloredverbatimCorrect, escapechar=@, caption=Solution for exercise 3]
use std::io;

fn main() {
  let @\hcb{mut}@ x: i32 = 10;
  x = 20;

  let y: f32 = @\hcb{2.5}@;

  println("X : ", x);
  println("Y : ", y);
}
\end{lstlisting}

\subsection*{Exercise 4}

Adding \token{10} to \token{250} exceeds the maximum value representable by a
\token{u8} (\token{255}), so the result overflows and wraps around:
\token{250 + 10 = 260}, and \token{260 mod 256 = 4}.

\begin{lstlisting}[style=bashVerb]
X : 4
\end{lstlisting}

\subsection*{Exercise 5}

\begin{lstlisting}[style=bashVerb]
A : 12
B : 12
C : 3.14
D : 3.14
E : 1000000
\end{lstlisting}

\token{a} carries no suffix, so it defaults to \token{i32}. \token{b} carries
the explicit suffix \token{u64}, so it is a 64-bit unsigned integer.
\token{c} carries no suffix, so it defaults to \token{f64}. \token{d} carries
the suffix \token{f}, so it is an \token{f32}. \token{e} is again an
\token{i32}; the underscores only group the digits for readability and do not
change the value or the type.

\subsection*{Exercise 6}

\begin{lstlisting}[style=bashVerb]
5
\end{lstlisting}

The word has four characters, but \token{é} is not part of the ASCII table
and requires two bytes to be encoded in UTF‑8, while every other character of
\textit{café} requires only one. The default string type is \token{[c8]},
whose \token{.len} reports a number of bytes, not a number of characters, so
the printed length is $3 + 2 = 5$.

\subsection*{Exercise 7}

\begin{lstlisting}[style=bashVerb]
Checking a
OK : true
\end{lstlisting}

Only \token{check("a", false)} is called. The \token{\&\&} operator is the
Boolean AND: as soon as its left operand is \token{false}, the whole
expression is necessarily \token{false}, so the right operand,
\token{check("b", true)}, is never evaluated, and \token{"Checking b"} is
never printed. Since the condition of the \token{if} is \token{false}, its
body is skipped and \token{ok} keeps its initial value \token{true}.
